\scanpage{3}
\mspage{1}
Complexe DR à p.d. et ombres des modèles.

% slanted
\margin{Plan conf. IHES en 1976}

1 Historique

a) Notion des formes diff. (\uline{Poincaré} ?) et formule de \uline{Stokes}
\[
\int_c d\omega=\int_{\partial c}\omega
\qquad\text{(d'où}\quad H^{\ast}_{\mathrm{DR}}(X)\longrightarrow H^{\ast}(X,\mathbb C)
\]
\margin{géom. diff\\ -- anal.\\ -- algébr.\\ \deleted{\unsure{top.}} arithm.\\ top. alg $\left\{\begin{array}{l}\mathrm{PL}\\1/2\,\unsure{\mathrm{simpl.}}\end{array}\right.$}

b) Th. de \uline{De Rham}. (Conjecturé par \uline{É. Cartan}). \unsure{Mais du point de présentation|Mais de présentation|Mais on perd la torsion}

Maintenant bien compris : th. des faisceaux, lemme Poincaré

c) Théorie de \uline{Hodge des} intégrales harmoniques (structure supplémentaire bigraduée sur $H^{\ast}_{\mathrm{DR}}(X)$ si $X$ kählérienne compacte)

d) Théorème de \uline{H. Cartan}-\uline{Serre} sur variétés de Stein (car lemme de Poincaré valable dans l'analytique)

e) Cas des variétés algébriques ou schémas sur \inserted{corps de} base (ou schéma de base général) : \uline{Dwork}, \uline{Washnitzer}-\uline{Monsky}, plus tard le yoga ``cristallin'' développé par \uline{Berthelot}, \uline{Illusie}, \deleted{\unsure{et repris par d'autres}}\inserted{\uline{Messing}, Mazur}
% slanted
\margin{cf aussi Hartshorne Herrera Ogus, Bloch (?)}

Ici on trouve des théories coh. qui ne sont plus à ``anneau de coeff''\unsure{\}} de car. nulle i.e contenant $\mathbb Q$ -- i.e. en principe les phénomènes de torsion ne sont pas perdus. \inserted{Du pt de vue Weil, \unsure{s'était|c'était} un défaut impardonnable !}

f) Th. de \uline{Grothendieck} pour variétés algébriques sur $\mathbb C$ (\deleted{\unsure{affiné}}\inserted{\uline{généralisé}} par \uline{Deligne}, \uline{Hartshorne} pour des coeff plus généraux) Ceci donne confiance (du moins en car. 0) en la signification topologique de la coh. de De Rham ``algébrique'' des schémas algébriques.

g) Complexe de De-Rham-Sullivan pour espaces topologiques généraux : \uline{formes différentielles singulières} $C^\infty$ \deleted{\illegible{2} $\unsure{C_{\mathrm{alg}}}$ \unsure{algébrique, resp.}} ($\mathbb C$-algébrique, resp. $\mathbb R$-algébrique, resp $\mathbb Q$-algébrique). Donne la coh. à coeff dans $\mathbb C$ (resp. $\mathbb R$, resp. $\mathbb Q$) (facile) Sullivan montre mieux.

\diagram{\begin{tikzcd}[ampersand replacement=\&,column sep=small,row sep=small]
C^\bullet_{\mathrm{DRS}\,\mathbb R\text{-}\mathrm{alg}} \arrow[r,phantom,"\subset"] \& C^\bullet_{\mathrm{DRS}\,C^\infty} \arrow[r,phantom,"\subset"] \arrow[d,phantom,"\supset"{rotate=-90}] \& C^\bullet_{\mathrm{DRS}} \\
\deleted{\unsure{C}} \& C^\bullet_{\mathrm{DR}} \&
\end{tikzcd}}
% archive sheet number: 2
